% EICL - Documento de preparación y línea base | LAB-000
% Revisión 0.3 (25/10/2026): actualizada con los resultados de LAB-001 y saneada
% para publicación según la norma de protección de la información (revisión 2, D-021).
% Compila con pdfLaTeX o XeLaTeX | Idioma: español
% Base: Documento Maestro de Transferencia de Conocimientos (v1.0, 18/10/2026),
%       LAB-000 v0.2 (23/10/2026) y memorias de LAB-001 (FortiGate, Huawei, Omada).
% ChkTeX: las advertencias 1, 8 y 13 generan falsos positivos con
% comandos válidos, códigos de activos, fechas y nombres como TP-Link.
% chktex-file 1
% chktex-file 8
% chktex-file 13
\documentclass[11pt,a4paper]{article}

\usepackage{iftex}
\ifPDFTeX
  \usepackage[utf8]{inputenc}
  \usepackage[T1]{fontenc}
  \IfFileExists{lmodern.sty}{\usepackage{lmodern}}{}
\else
  \usepackage{fontspec}
\fi
% Español: con babel-spanish instalado se usan las opciones de siempre;
% si falta, babel carga el idioma desde sus archivos .ini.
\IfFileExists{spanish.ldf}
  {\usepackage[spanish,es-nodecimaldot,es-noshorthands]{babel}}
  {\usepackage{babel}\babelprovide[import,main]{spanish}}
\usepackage[margin=2.3cm,headheight=18pt]{geometry}
\usepackage[table]{xcolor}
\usepackage{booktabs,tabularx,longtable,array}
\usepackage{graphicx}
\usepackage{hyperref}
\usepackage{fancyhdr}
\usepackage{titlesec}
\usepackage{enumitem}
\usepackage{lastpage}
\usepackage{listings}
\ifPDFTeX\IfFileExists{lmodern.sty}{\usepackage{microtype}}{\usepackage[expansion=false]{microtype}}\else\usepackage{microtype}\fi
\usepackage{ragged2e}

\definecolor{navy}{HTML}{14283F}
\definecolor{blue}{HTML}{16608B}
\definecolor{muted}{HTML}{5D6977}
\definecolor{light}{HTML}{EAF0F4}
\definecolor{accent}{HTML}{16A3A1}
\definecolor{okgreen}{HTML}{2E7D32}
\definecolor{amber}{HTML}{B26A00}
\definecolor{alertred}{HTML}{B3261E}
\definecolor{violet}{HTML}{6A4C93}

\hypersetup{
  colorlinks=true,
  linkcolor=navy,
  urlcolor=blue,
  pdftitle={EICL -- LAB-000 Preparación, inventario y línea base},
  pdfauthor={Carlos Eduardo Abundis Lira}
}

\titleformat{\section}{\Large\bfseries\color{navy}}{\thesection}{0.8em}{}
\titleformat{\subsection}{\large\bfseries\color{blue}}{\thesubsection}{0.8em}{}
\setlist[itemize]{leftmargin=1.6em,itemsep=2pt,topsep=4pt}
\setlist[enumerate]{leftmargin=1.8em,itemsep=2pt,topsep=4pt}
\setlength{\parindent}{0pt}
\setlength{\parskip}{6pt}
\renewcommand{\arraystretch}{1.25}

\lstset{
  basicstyle=\ttfamily\small,
  frame=single,
  rulecolor=\color{light},
  backgroundcolor=\color{light!40},
  columns=fullflexible,
  keepspaces=true,
  xleftmargin=4pt,
  xrightmargin=4pt
}

\pagestyle{fancy}
\fancyhf{}
\fancyhead[L]{\footnotesize\textcolor{muted}{EICL | LAB-000}}
\fancyhead[R]{\footnotesize\textcolor{muted}{Preparación y línea base del homelab}}
\fancyfoot[L]{\footnotesize\textcolor{muted}{Rev.\ 0.3 | Etapa 00 | Información saneada}}
\fancyfoot[R]{\footnotesize\textcolor{muted}{Pág.\ \thepage\ de \pageref{LastPage}}}
\renewcommand{\headrulewidth}{0.3pt}

% --- Marcadores ---
\newcommand{\pc}{\textcolor{muted}{\textit{[POR CONFIRMAR]}}}
% --- Estados normalizados (Documento maestro, sección 2.1 y 5.1) ---
\newcommand{\estadobase}[2]{\textcolor{#1}{\textbf{\footnotesize #2}}}
\newcommand{\disp}{\estadobase{okgreen}{DISPONIBLE}}
\newcommand{\adq}{\estadobase{blue}{ADQUIRIDO}}
\newcommand{\porval}{\estadobase{amber}{POR VALIDAR}}
\newcommand{\prop}{\estadobase{violet}{PROPUESTO}}
\newcommand{\ant}{\estadobase{muted}{INSTALADO ANTERIORMENTE}}
\newcommand{\usant}{\estadobase{muted}{UTILIZADO ANTERIORMENTE}}
\newcommand{\pend}{\estadobase{alertred}{PENDIENTE}}
\newcommand{\curso}{\estadobase{amber}{EN CURSO}}
\newcommand{\compl}{\estadobase{okgreen}{COMPLETADA}}
\newcommand{\histo}{\estadobase{muted}{HISTÓRICO}}

% --- Columnas ---
\newcolumntype{L}[1]{>{\RaggedRight\arraybackslash}p{#1}}
\newcolumntype{Z}{>{\RaggedRight\arraybackslash}X}

% --- Recuadros ---
\newcommand{\notebox}[2]{%
  \par\noindent\fcolorbox{blue}{light}{%
    \parbox{\dimexpr\linewidth-2\fboxsep-2\fboxrule\relax}{%
      {\small\bfseries\color{blue}\MakeUppercase{#1}}\\[2pt]\small #2}}\par}
\newcommand{\alertbox}[2]{%
  \par\noindent\fcolorbox{alertred}{alertred!6}{%
    \parbox{\dimexpr\linewidth-2\fboxsep-2\fboxrule\relax}{%
      {\small\bfseries\color{alertred}\MakeUppercase{#1}}\\[2pt]\small #2}}\par}

\begin{document}
\hypersetup{pageanchor=false}

\begin{titlepage}
  \thispagestyle{empty}
  \vspace*{2.0cm}
  {\color{accent}\rule{2cm}{5pt}}\\[1.0cm]
  {\sffamily\Large\bfseries\color{blue}
    ENTERPRISE INFRASTRUCTURE \& CYBERSECURITY}\\[0.08cm]
  {\sffamily\Large\bfseries\color{blue}
    LABORATORY}\\[0.55cm]
  {\sffamily\fontsize{29}{34}\selectfont\bfseries\color{navy}
    LAB-000\\ Preparación, inventario\\ y línea base}\\[0.7cm]
  {\Large\color{muted}Etapa 00 | Recepción, validación y reconstrucción}\\[1.5cm]

  \renewcommand{\arraystretch}{1.6}
  \begin{tabularx}{\textwidth}{>{\bfseries}p{3.4cm}X}
    Código & EICL-LAB-000 \\
    Versión & 0.3 -- Línea base actualizada con LAB-001 y saneada \\
    Fecha & 25 de octubre de 2026 \\
    Responsable & Carlos Eduardo Abundis Lira \\
    Proyecto & EICL-LAB | Horizonte inicial de 90 días \\
    Estado & Etapa 00 en curso; en pausa hasta cerrar LAB-001 (D-012) \\
    Clasificación & Portafolio técnico | Información saneada \\
  \end{tabularx}

  \vfill
  {\color{navy}\rule{\textwidth}{1pt}}\\[0.25cm]
  \centering {\small\color{muted}
    Redes \quad / \quad Virtualización \quad / \quad Seguridad \quad / \quad Observabilidad}
\end{titlepage}
\hypersetup{pageanchor=true}

\tableofcontents
\clearpage

% =====================================================================
\section{Control documental}
\begin{tabularx}{\textwidth}{L{1.7cm} L{2.2cm} Z}
  \toprule
  \rowcolor{light}\textbf{Revisión} & \textbf{Fecha} & \textbf{Cambio}\\
  \midrule
  0.1 & 2026-10-14 & Alta del formato como LAB-001 (levantamiento AS-IS); datos iniciales sujetos a comprobación.\\
  0.2 & 2026-10-23 & Reetiquetado como LAB-000 (Etapa 00). Alineado al documento maestro v1.0: tres nodos OptiPlex, FortiGate y Huawei como equipos adquiridos, rack y cableado, condición de reconstrucción desde cero, plan VLAN propuesto, riesgos y criterios de cierre.\\
  0.3 & 2026-10-25 & Incorpora los resultados de LAB-001: FortiGate 50E, Huawei S5735 y Omada ES208G recibidos, levantados y configurados; versiones de software comprobadas (actividades 6 y 7, VAL-004 y VAL-005); red doméstica documentada (VAL-007); plan VLAN/IP implementado; matriz de enlaces con los puertos reales. Corrige el papel de los routers domésticos (el Archer C3200 es el router del cuarto). Saneada para publicación (D-021): sin nombre ni modelo del equipo del proveedor de internet, versiones de firmware generalizadas y tabla «Protección de la información».\\
  1.0 & Pendiente & Etapa 00 cerrada con inventario físico validado y arquitectura inicial aprobada.\\
  \bottomrule
\end{tabularx}

\subsection{Documentos de referencia}
\begin{itemize}
  \item Documento Maestro de Transferencia de Conocimientos del proyecto EICL-LAB, versión 1.0 del 18/10/2026. Línea base del proyecto.
  \item EICL-LAB-001 Levantamiento y arquitectura de infraestructura, v0.1 (14/10/2026). Sustituido por este documento para la Etapa 00.
  \item Memorias técnicas de LAB-001, evidencia de los datos de red de esta revisión:
  \begin{itemize}
    \item Componente A: \texttt{\small EICL-LAB-001-fortigate-levantamiento.pdf}
    \item Componente B: \texttt{\small EICL-LAB-001-huawei-levantamiento.pdf}
    \item Componente C: \texttt{\small EICL-LAB-001-omada-levantamiento.pdf}
  \end{itemize}
\end{itemize}

\subsection{Estados utilizados}
\begin{tabularx}{\textwidth}{L{4.6cm} Z}
  \toprule
  \rowcolor{light}\textbf{Estado} & \textbf{Significado}\\
  \midrule
  \disp & Se cuenta físicamente con el equipo.\\
  \adq & Compra realizada; entrega o recepción pendiente de verificar.\\
  \porval & Existe registro del equipo; falta verificar funcionamiento o especificaciones.\\
  \prop & Contemplado como posible adquisición o implementación.\\
  \ant & Tecnología utilizada antes de la reconstrucción.\\
  \pend & Todavía no implementado en el nuevo laboratorio.\\
  \curso & Hay actividades realizadas, pero faltan criterios de cierre.\\
  \compl & Actividad o validación cerrada con evidencia.\\
  \bottomrule
\end{tabularx}

% =====================================================================
\section{Resumen ejecutivo y alcance}
\textbf{Objetivo de LAB-000.} Recibir, inspeccionar e inventariar el equipo del laboratorio, documentar riesgos y dejar lista la arquitectura inicial para implementar LAB-001. Es la Etapa 00 del roadmap de 90 días.

\textbf{Condición de inicio.} El laboratorio se reconstruye desde cero: se eliminan configuraciones experimentales, se reinstalan plataformas, se rediseñan direccionamiento y VLAN, y se documenta cada componente desde su instalación. Los servicios usados antes (Proxmox VE, OPNsense, Wazuh, Tailscale) \textbf{no} cuentan como implementados hasta reinstalarse y validarse.

\notebox{Situación al 25/10/2026}{La parte de red de la Etapa 00 avanzó junto con LAB-001: el firewall, el switch principal y el switch de acceso secundario están recibidos, levantados, configurados y validados de punta a punta. LAB-000 queda en pausa hasta cerrar LAB-001 (D-012); siguen pendientes los nodos OptiPlex, el material del cable y el montaje definitivo en rack.}

\newpage
\textbf{Alcance.}
\begin{itemize}
  \item Tres nodos Dell OptiPlex 5070 Micro y componente de reparación (placa base).
  \item FortiGate 50E, Huawei S5735-L8P4S-A1, TP-Link Omada ES208G y TP-Link TL-SG108.
  \item Equipos de la red doméstica que delimitan el laboratorio (módem del proveedor de internet, router del cuarto y AP adicionales).
  \item Rack de 12U, charolas, cableado, conectores y protección eléctrica.
  \item Respaldo selectivo previo al borrado, roles de nodos y plan de direccionamiento.
\end{itemize}

\textbf{Fuera de alcance.} Configuración de VLAN, políticas de firewall y despliegue de hipervisores (corresponden a LAB-001 y LAB-002); producción empresarial externa; configuraciones o secretos de terceros; cambios no evaluados en la red doméstica.

% =====================================================================
\section{Inventario de infraestructura física}
\subsection{Equipo de cómputo}
{\small
\begin{longtable}{L{1.5cm} L{3.0cm} L{3.9cm} L{2.8cm} L{3.0cm}}
  \toprule
  \rowcolor{light}\textbf{ID} & \textbf{Equipo} & \textbf{Especificaciones registradas} & \textbf{Estado} & \textbf{Rol previsto}\\
  \midrule
  \endhead
  SRV-01 & Dell OptiPlex 5070 Micro & Core i5 9.ª gen., 16\,GB RAM, SSD 256\,GB & \disp; funcionamiento \pc & Nodo 01 -- Proxmox VE\\
  SRV-02 & Dell OptiPlex 5070 Micro & Core i5 9.ª gen., 16\,GB RAM, SSD 256\,GB & \disp; funcionamiento \pc & Nodo 02 -- VMware ESXi\\
  SRV-03 & Dell OptiPlex 5070 Micro & i5 9.ª gen., 16\,GB RAM, SSD 256\,GB según configuración prevista & \adq; recepción y especificaciones reales \pc & Nodo 03 -- flexible\\
  REP-01 & Placa base compatible OptiPlex 5070 Micro & Referencia P01GV & \adq; recepción e instalación \pc & Reparación de hardware\\
  \bottomrule
\end{longtable}
}

Dimensiones de referencia de cada OptiPlex Micro: 182 $\times$ 178 $\times$ 36\,mm.

\textbf{Observaciones.}
\begin{itemize}
  \item Existe antecedente de falla de una placa base; no está registrado cuál de los nodos la presentó.
  \item Verificar cuántos nodos están realmente operativos. No asumir que los tres están listos.
  \item Inventariar físicamente: procesador exacto, módulos de RAM, SSD, NIC, capacidad de expansión y versión de BIOS.
  \item Los puertos GE0/0/2 a GE0/0/4 del switch principal están reservados y apagados para SRV-01 a SRV-03 hasta LAB-002. Los servidores quedan fuera del rack hasta revisar su POST.
\end{itemize}

\newpage
\subsection{Networking}
{\small
\begin{longtable}{L{1.5cm} L{3.0cm} L{3.9cm} L{2.8cm} L{3.0cm}}
  \toprule
  \rowcolor{light}\textbf{ID} & \textbf{Equipo} & \textbf{Especificaciones registradas} & \textbf{Estado} & \textbf{Función}\\
  \midrule
  \endhead
  FW-01 & Fortinet FortiGate 50E & FortiOS 6.x comprobado tras reset de fábrica (memoria LAB-001-A); sin FortiGuard; licencia y registro \pc & \disp; configurado y validado (FGT-LAB-01) & Firewall perimetral del laboratorio y gateway de las VLAN\\
  SW-01 & Huawei S5735-L8P4S-A1 & 8 GE PoE+ y 4 SFP; 300 $\times$ 220 $\times$ 43.6\,mm; VRP V200R020 comprobado (memoria LAB-001-B); PoE+ confirmado & \disp; configurado y validado (SW-LAB-01) & Switch principal administrable\\
  SW-02 & TP-Link Omada ES208G & Easy Managed, 8 puertos Gigabit; hardware v1.6 (equivale a V1.0), firmware 1.0.x (memoria LAB-001-C) & \disp; configurado y validado (SW-LAB-02) & Switch de acceso secundario\\
  SW-03 & TP-Link TL-SG108 & 8 puertos Gigabit, no administrable; modelo exacto \pc & \disp; fuera de LAB-001 por ahora & Switching simple\\
  RTR-01 & TP-Link Archer C3200 & Router/AP Wi-Fi; red del cuarto 192.168.1.0/24 & \disp; en uso & Router del cuarto; su LAN2 conecta la WAN del FortiGate\\
  RTR-02 & TP-Link Archer GE800 & Router/AP Wi-Fi & \disp & Red doméstica; papel actual \pc\\
  AP-01 & TP-Link Archer C5200 & Usado antes como router/AP & \disp & WLAN\\
  ISP-01 & Equipo del proveedor de internet & Modelo omitido en la publicación & \disp & WAN / acceso a internet (192.168.0.0/24)\\
  \bottomrule
\end{longtable}
}

\notebox{Resultado de las validaciones de networking}{\textbf{FortiGate:} versión de FortiOS comprobada; sin FortiGuard, así que los laboratorios no dependen de servicios UTM con licencia. Registro en FortiCare \pc. \textbf{Huawei:} versión VRP comprobada; llegó sin configuración guardada y sin daños físicos (reportado). \textbf{Omada:} la etiqueta confirma el modelo ES208G (la denominación DS108GE del inventario anterior queda descartada); soporta 802.1Q, Management VLAN, LAG, Port Mirroring, Cable Test y Loop Prevention, pero no STP, SNMP, syslog ni CLI.}

\subsection{Rack, cableado y energía}
{\small
\begin{longtable}{L{3.0cm} L{5.2cm} L{2.8cm} L{3.7cm}}
  \toprule
  \rowcolor{light}\textbf{Elemento} & \textbf{Especificaciones} & \textbf{Estado} & \textbf{Función}\\
  \midrule
  \endhead
  Rack genérico & Rack abierto 12U, blanco & \porval; en uso en el componente D de LAB-001 (reportado) & Alojamiento de equipos\\
  Charolas & 4 charolas para rack & \adq; recepción \pc & Montaje de equipos compactos\\
  Cable Ethernet plano & Cable UTP; material \pc & \adq & Cableado físico; no usado en enlaces del laboratorio hasta verificar el material\\
  Patch cords de fábrica & Existencias \pc & \disp & Troncales del laboratorio (FortiGate--Huawei--Omada)\\
  Conectores RJ45 & Cantidad y tipo \pc & \porval & Terminación de cables\\
  Regulador / UPS & Modelo y disponibilidad \pc & \porval & Protección eléctrica\\
  PDU / multicontactos & Disponibilidad y capacidad 12 puertos / 15A & \adq & Distribución eléctrica\\
  \bottomrule
\end{longtable}
}

\alertbox{Cableado y PoE}{Al utilizar cableado de tipo CCA (aluminio revestido de cobre), no se usará en enlaces PoE. PoE queda apagado en todos los puertos del switch Huawei (D-009); si algún día se necesita, se activa puerto por puerto y con cable verificado.}

\textbf{Restricciones físicas:} confirmar medidas internas útiles del rack; validar ventilación sin bloquear entradas de aire; verificar peso y tamaño soportados por las charolas; identificar con etiqueta cada equipo y cable.

% =====================================================================
\section{Plataformas y servicios: estado de reconstrucción}
{\small
\begin{longtable}{L{3.2cm} L{4.3cm} L{3.8cm} L{3.4cm}}
  \toprule
  \rowcolor{light}\textbf{Software} & \textbf{Uso} & \textbf{Estado} & \textbf{Requisito}\\
  \midrule
  \endhead
  Proxmox VE & Hipervisor principal (Nodo 01) & \ant & Reinstalación pendiente (LAB-002)\\
  Proxmox VE 9.2.2 & Versión de la configuración anterior & \histo & Versión final \pc\\
  VMware ESXi & Segundo hipervisor (Nodo 02) & \prop & Compatibilidad y licencia \pc\\
  Wazuh & SIEM / monitoreo de seguridad & \usant & Recursos suficientes (LAB-006)\\
  Tailscale & Acceso remoto seguro & \usant & Escenario por definir (LAB-010)\\
  Debian Linux & Servidores y servicios & \prop & Obtener e instalar\\
  Red Hat Enterprise Linux & Administración Linux empresarial & \prop & Suscripción Developer \pc\\
  Docker & Contenedores & \prop & Definir host\\
  Proxmox Backup Server & Respaldos & \prop & Destino de respaldo \pc\\
  \bottomrule
\end{longtable}
}

No dar por habilitados respaldos, alta disponibilidad ni registros centralizados hasta comprobarlos en el nuevo laboratorio.
\newpage
% =====================================================================
\section{Arquitectura física objetivo (TO-BE)}
La siguiente topología combina lo ya implementado en LAB-001 (red) con lo previsto (nodos). Los nodos \textbf{no} están instalados.

\begin{lstlisting}
Internet
   |
Equipo del ISP (ISP-01) .......... red de casa 192.168.0.0/24
   |
Archer C3200 (RTR-01) ............ red del cuarto 192.168.1.0/24
   |  LAN2
[wan1] FortiGate 50E (FW-01) ...... implementado
   |  lan5: trunk 802.1Q
[GE0/0/1] Huawei S5735 (SW-01) .... implementado
   |-- GE0/0/2  Nodo 01: Proxmox VE (SRV-01) ....... previsto (LAB-002)
   |-- GE0/0/3  Nodo 02: VMware ESXi (SRV-02) ...... previsto (LAB-002)
   |-- GE0/0/4  Nodo 03: plataforma flexible ....... previsto (LAB-002)
   |-- GE0/0/5  Omada ES208G (SW-02), puerto 1 .... implementado
   |               `-- puertos 2-5: acceso MGMT, SRV y LAB
   `-- GE0/0/6-8  acceso MGMT, SRV y LAB
\end{lstlisting}

\begin{itemize}
  \item \textbf{FortiGate:} gateway de todas las VLAN y punto de control entre segmentos (router-on-a-stick, D-002).
  \item \textbf{Huawei:} switch principal administrable, en capa 2 (D-006).
  \item \textbf{Omada:} switch de acceso secundario en modo standalone (D-013).
\end{itemize}

\subsection{Distribución lógica de nodos}
\begin{tabularx}{\textwidth}{L{2.4cm} L{3.4cm} Z}
  \toprule
  \rowcolor{light}\textbf{Nodo} & \textbf{Plataforma} & \textbf{Uso y condiciones}\\
  \midrule
  Nodo 01 & Proxmox VE & Servicios virtualizados, VMs Linux, infraestructura de red y aplicaciones.\\
  Nodo 02 & VMware ESXi & VMs de práctica y redes virtuales. Sujeto a compatibilidad de CPU, NIC y almacenamiento, y a licencia. Sin vCenter, vMotion, DRS ni HA empresarial.\\
  Nodo 03 & Sin asignación definitiva & Opciones: Debian, RHEL, Docker, servidor de respaldos, servicios de infraestructura o segundo nodo Proxmox para clúster y HA.\\
  \bottomrule
\end{tabularx}

\subsection{Matriz de enlaces}
\begin{tabularx}{\textwidth}{L{1.9cm} Z Z Z}
  \toprule
  \rowcolor{light}\textbf{ID} & \textbf{Origen / puerto} & \textbf{Destino / puerto} & \textbf{Observación}\\
  \midrule
  LINK-001 & RTR-01 Archer C3200 / LAN2 & FW-01 / wan1 & WAN del laboratorio hacia la red del cuarto (cable C-02)\\
  LINK-002 & FW-01 / lan5 & SW-01 / GE0/0/1 & Trunk 802.1Q, VLAN 10--50, nativa 999; 1000 Mbps full (cable C-03)\\
  LINK-003 & SW-01 / GE0/0/2 & SRV-01 / NIC & Reservado y apagado hasta LAB-002; velocidad \pc\\
  LINK-004 & SW-01 / GE0/0/3 & SRV-02 / NIC & Reservado y apagado hasta LAB-002; velocidad \pc\\
  LINK-005 & SW-01 / GE0/0/4 & SRV-03 / NIC & Reservado y apagado hasta LAB-002; velocidad \pc\\
  LINK-006 & SW-01 / GE0/0/5 & SW-02 / puerto 1 & Trunk 802.1Q, VLAN 10, 20 y 50, nativa 999 (cable C-04); velocidad \pc\\
  \bottomrule
\end{tabularx}

Los puertos se trazaron físicamente durante LAB-001.\\
Tabla completa: \texttt{docs/port-mapping.md} del repositorio.

% =====================================================================
\section{Arquitectura lógica}
\subsection{Red doméstica actual}
\begin{tabularx}{\textwidth}{L{3.0cm} L{3.0cm} L{2.0cm} Z}
  \toprule
  \rowcolor{light}\textbf{Segmento} & \textbf{CIDR} & \textbf{VLAN} & \textbf{Gateway / estado}\\
  \midrule
  Red de casa & 192.168.0.0/24 & No aplica & Equipo del ISP en 192.168.0.1 (VAL-007)\\
  Red del cuarto & 192.168.1.0/24 & No aplica & Archer C3200 en 192.168.1.1; DHCP .100--.199; WAN del FortiGate fija en 192.168.1.2 (VAL-007)\\
  \bottomrule
\end{tabularx}

\subsection{Plan de VLAN (implementado en LAB-001)}
\begin{tabularx}{\textwidth}{L{1.4cm} Z L{2.6cm} L{2.2cm} L{2.6cm}}
  \toprule
  \rowcolor{light}\textbf{VLAN} & \textbf{Uso} & \textbf{Subred} & \textbf{Gateway} & \textbf{DHCP}\\
  \midrule
  10 & Administración (MGMT) & 10.10.10.0/24 & 10.10.10.1 & .100--.120\\
  20 & Servidores (SRV) & 10.10.20.0/24 & 10.10.20.1 & .100--.199\\
  30 & Monitoreo / SOC (MON) & 10.10.30.0/24 & 10.10.30.1 & Sin DHCP\\
  40 & DMZ & 10.10.40.0/24 & 10.10.40.1 & Sin DHCP\\
  50 & Laboratorio y clientes (LAB) & 10.10.50.0/24 & 10.10.50.1 & .100--.199\\
  999 & Nativa de trunks y puertos sin uso & Sin IP & No aplica & No aplica\\
  60 / 70 & Respaldos (LAB-008) / clúster Proxmox (LAB-012) & Reservadas & -- & --\\
  \bottomrule
\end{tabularx}

\notebox{Estado del plan}{Implementado y validado de punta a punta en LAB-001 (V1 a V8). Rango 10.10.0.0/16 con el tercer octeto igual al ID de VLAN (D-001); dentro de cada /24: .1 gateway, .2--.9 equipos de red, .10--.49 servidores e hipervisores, .100--.199 DHCP. Dominio interno \texttt{eicl.internal}. Pendiente para esta etapa: roles definitivos de los nodos.}

% =====================================================================
\section{Actividades de la Etapa 00}
{\small
\begin{longtable}{L{0.7cm} L{6.4cm} L{2.7cm} L{4.7cm}}
  \toprule
  \rowcolor{light}\textbf{\#} & \textbf{Actividad} & \textbf{Estado} & \textbf{Evidencia}\\
  \midrule
  \endhead
  1 & Recibir equipos pendientes (SRV-03, REP-01, FW-01, SW-01, rack y charolas) & \curso & FW-01 y SW-01 recibidos; el resto \pc\\
  2 & Inspeccionar hardware (daños, puertos, accesorios, fuentes) & \curso & FW-01 y SW-01 inspeccionados (reportado); nodos pendientes\\
  3 & Verificar arranque y POST de cada OptiPlex & \pend & Sin registrar\\
  4 & Revisar la reparación de la placa base & \pend & Sin registrar\\
  5 & Identificar CPU, RAM, SSD, NIC y BIOS por nodo & \pend & Sin registrar\\
  6 & Comprobar firmware del FortiGate & \compl & Memoria LAB-001-A, E-01\\
  7 & Comprobar versión VRP del Huawei & \compl & Memoria LAB-001-B, EVID-002\\
  8 & Verificar material del cable y conectores & \pend & Se usan patch cords de fábrica (RSK-005)\\
  9 & Organizar montaje y ventilación en el rack & \pend & Componente D de LAB-001 en curso\\
  10 & Crear inventario (este documento como base) & \curso & LAB-000 v0.3\\
  11 & Guardar respaldos selectivos de datos relevantes & \pend & La configuración previa del FortiGate se descartó a propósito (era de terceros)\\
  12 & Preparar la eliminación de configuraciones anteriores & \curso & FortiGate: reset hecho; Huawei y Omada llegaron sin configuración; nodos y TL-SG108 pendientes\\
  13 & Definir direccionamiento y roles & \curso & Plan VLAN/IP implementado; roles de nodos pendientes\\
  14 & Crear estructura inicial de documentación & \curso & Plantilla EICL, memorias de LAB-001 y repositorio público\\
  \bottomrule
\end{longtable}
}

\subsection{Respaldo previo al borrado}
Antes de eliminar configuraciones experimentales, revisar y conservar lo que tenga valor:
\begin{itemize}
  \item Exportar inventario de VM/CT y configuración de red de Proxmox anterior, si los discos siguen accesibles.
  \item Exportar configuración de OPNsense (XML) y de los switches TP-Link, solo como referencia.
  \item Rescatar notas, diagramas y capturas del laboratorio anterior.
  \item Guardar todo fuera de los nodos que se reinstalarán; nunca publicar secretos contenidos en esos archivos.
\end{itemize}

\alertbox{Acción destructiva}{La reinstalación borra los discos de los nodos. Confirmar que el respaldo selectivo está completo y legible antes de iniciar cualquier borrado.}

\newpage
% =====================================================================
\section{Plan de validación}
{\small
\begin{longtable}{L{1.6cm} L{4.2cm} L{5.6cm} L{3.3cm}}
  \toprule
  \rowcolor{light}\textbf{Prueba} & \textbf{Objetivo} & \textbf{Evidencia esperada u obtenida} & \textbf{Estado}\\
  \midrule
  \endhead
  VAL-001 & Arranque de SRV-01, 02 y 03 & Foto o captura del POST y pantalla de BIOS con fecha & \pend\\
  VAL-002 & Especificaciones por nodo & CPU, RAM, SSD, NIC y versión de BIOS registradas en \texttt{INVENTORY.md} & \pend\\
  VAL-003 & Placa base reparada & Nodo afectado arranca y reconoce RAM, SSD y NIC & \pend\\
  VAL-004 & Versión FortiOS y registro & \texttt{get system status} saneada (memoria LAB-001-A) & \compl{} (versión); registro \pc\\
  VAL-005 & Versión VRP del Huawei & \texttt{display version} saneada (memoria LAB-001-B) & \compl\\
  VAL-006 & Material del cable & Prueba o especificación que distinga cobre de CCA & \pend\\
  VAL-007 & Red doméstica & CIDR, gateway y DHCP de las dos redes domésticas (sección 6.1) & \compl\\
  \bottomrule
\end{longtable}
}

Las pruebas de conectividad, DNS y segmentación pasaron a LAB-001 y están cerradas allí (V1 a V8).

% =====================================================================
\section{Riesgos y bloqueos identificados}
{\small
\begin{longtable}{L{1.5cm} L{2.3cm} L{5.4cm} L{5.5cm}}
  \toprule
  \rowcolor{light}\textbf{ID} & \textbf{Tipo} & \textbf{Riesgo} & \textbf{Mitigación}\\
  \midrule
  \endhead
  RSK-001 & Hardware & Falla previa de placa base; no está confirmado cuántos nodos operan & Ejecutar VAL-001 y VAL-003 antes de asignar roles\\
  RSK-002 & Licencia & ESXi puede no ser compatible o licenciable en el Nodo 02 & Validar CPU, NIC y licencia; alternativa: Nodo 02 como segundo Proxmox\\
  RSK-003 & Firmware & Funciones UTM sin licencia (sin FortiGuard) & Mitigado: versión FortiOS confirmada; se planea solo con funciones base\\
  RSK-004 & Firmware & Versión VRP del Huawei desconocida & Cerrado: versión confirmada (VAL-005)\\
  RSK-005 & Hardware & Cable plano posiblemente CCA, inadecuado para PoE & Ejecutar VAL-006; PoE apagado (D-009); patch cords de fábrica en los troncales\\
  RSK-006 & Seguridad & Interferencia del laboratorio con la red doméstica & Mitigado: FortiGate como frontera y política 1 que bloquea las redes domésticas, validada (V2)\\
  RSK-007 & Configuración & Pérdida de datos útiles al reconstruir & Respaldo selectivo verificado antes del borrado\\
  RSK-008 & Recursos & 16\,GB por nodo limitan servicios pesados (Wazuh, Zabbix) & Medir consumo antes de LAB-005 y LAB-006\\
  RSK-009 & Tiempo / presupuesto & Horas semanales y presupuesto mensual sin definir & Confirmar con el propietario \pc\\
  \bottomrule
\end{longtable}
}

\newpage
% =====================================================================
\section{Criterios de cierre de LAB-000}
LAB-000 pasa a COMPLETADA solo cuando exista evidencia de:
\begin{enumerate}
  \item Todos los equipos recibidos e inspeccionados, con estado actualizado en el inventario.
  \item Número de nodos operativos confirmado y especificaciones reales registradas.
  \item Versiones de FortiOS y VRP documentadas. \textit{(Cumplido.)}
  \item Material del cableado identificado y decisión sobre PoE registrada. \textit{(Decisión sobre PoE registrada: D-009; material pendiente.)}
  \item Respaldo selectivo completado antes de cualquier borrado.
  \item Roles de nodos y plan de VLAN/IP aprobados. \textit{(Plan VLAN/IP implementado; roles pendientes.)}
  \item Riesgos documentados y arquitectura inicial lista para LAB-001. \textit{(Cumplido.)}
\end{enumerate}

\subsection{Estado de seguimiento}
\begin{tabularx}{\textwidth}{L{4.2cm} Z}
  \toprule
  \rowcolor{light}\textbf{Campo} & \textbf{Contenido}\\
  \midrule
  Proyecto actual & LAB-000 -- Etapa 00 (en pausa hasta cerrar LAB-001, D-012)\\
  Estado & \curso\\
  Avance conseguido & Firewall, switch principal y switch de acceso recibidos, levantados, configurados y validados (LAB-001); red doméstica documentada; plan VLAN/IP implementado\\
  Evidencia obtenida & Memorias de LAB-001 (componentes A, B y C)\\
  Pendientes & Actividades 1--5, 8, 9, 11 y 13 (roles); VAL-001, VAL-002, VAL-003 y VAL-006\\
  Bloqueos & Ninguno confirmado; ver riesgos RSK-001, RSK-002, RSK-005, RSK-007 a RSK-009\\
  Siguiente paso concreto & Al cerrar LAB-001: arrancar cada OptiPlex fuera del rack, registrar POST y BIOS (VAL-001) y sus especificaciones reales (VAL-002)\\
  Criterio de éxito & Número de nodos operativos y especificaciones registradas en \texttt{INVENTORY.md} con fecha y evidencia saneada\\
  \bottomrule
\end{tabularx}

% =====================================================================
\section{Protección de la información}
Este documento es la línea base de la Etapa 00. Los datos marcados como \pc{} no deben presentarse como implementación validada. Se trata de un laboratorio personal, no de una implementación empresarial.

Criterio del proyecto: se censura cualquier dato que no sea de dominio público o que pueda representar una vulnerabilidad, aunque sea mínima. \textbf{Canales revisados en esta revisión (0.3):} texto general (incluidas tablas y metadatos del PDF) y bloques de texto tipo consola; el documento no contiene capturas.

{\small
\begin{longtable}{L{3.4cm} L{6.0cm} L{5.6cm}}
  \toprule
  \rowcolor{light}\textbf{Dato} & \textbf{Tratamiento} & \textbf{Motivo de protección}\\
  \midrule
  \endhead
  Nombre comercial del proveedor de internet y modelo de su equipo & Omitidos (presentes en la v0.2, que no se publica); se usa «equipo del proveedor de internet». & Ubican la red doméstica y orientan ataques a un equipo concreto.\\
  Versiones exactas de firmware de los equipos de red & Generalizadas («FortiOS 6.x», «VRP V200R020», «firmware 1.0.x»). & Junto con la topología, orientan la búsqueda de vulnerabilidades públicas (regla D-022 de la norma de protección).\\
  Contraseñas, tokens, llaves, números de serie y respaldos & No se incluyen. & Acceso directo o identificación de los equipos.\\
  IP de la PC de administración & No aparece. & Es el origen autorizado para administrar el laboratorio.\\
  Configuraciones de terceros & No se incluyen (la configuración previa del FortiGate se descartó sin respaldar). & No son propias.\\
  \midrule
  \multicolumn{3}{l}{\textbf{Visible a propósito}}\\
  \midrule
  Plan de direccionamiento y redes domésticas genéricas (192.168.0.0/24, 192.168.1.0/24) & Secciones 5 y 6. & Diseño documentado y rangos privados de fábrica.\\
  Modelos de los equipos propios (OptiPlex, FortiGate, Huawei, Omada, Archer) & Inventario. & Inventario técnico del laboratorio; ninguno se expone a internet.\\
  Herramientas usadas antes (Proxmox VE 9.2.2, Tailscale, Wazuh, OPNsense) & Sección 4. & Historial de habilidades; ninguna está instalada en el laboratorio actual.\\
  Nombre del autor & Portada. & Autoría del portafolio.\\
  \bottomrule
\end{longtable}
}

\end{document}
